% ---------- Datei: chapters/01_einleitung.tex ----------
\chapter{Einleitung}

\section{Projektumfeld}
Das Projekt wird im Rahmen meiner Ausbildung zur Fachinformatikerin für Systemintegration bei der Söldner Consult GmbH in Nürnberg durchgeführt. Das Unternehmen ist ein spezialisiertes IT-Beratungshaus mit Fokus auf Cloud Computing (GCP, Azure, AWS), Datacenter und Networking.

Der Kunde ist ein mittelständisches Unternehmen aus dem Finanzsektor. Aufgrund der Branchenzugehörigkeit und strenger Compliance-Richtlinien unterliegt die IT-Infrastruktur hohen Sicherheitsvorgaben. Betriebsgeheimnisse, exakte Systemnamen und reale IP-Adressen des Kunden wurden in dieser Dokumentation aus Datenschutzgründen anonymisiert.

\section{Projektbeschreibung und Zielsetzung}
Die aktuelle Netzwerkarchitektur des Kunden in der Google Cloud Platform (GCP) ist historisch gewachsen. Sie besteht aus mehreren isolierten Virtual Private Clouds (VPCs) und dezentral konfigurierten VPN-Verbindungen zu den On-Premises-Standorten. Diese Struktur ist wartungsintensiv, fehleranfällig und skaliert schlecht bei der Anbindung neuer Zweigstellen.

Ziel dieses Projekts ist die Konzeption und Implementierung einer zentral verwalteten Hybrid-Cloud-Netzwerkarchitektur. Als technologische Basis dient das Google Cloud Network Connectivity Center (NCC). Durch den Aufbau einer Hub-and-Spoke-Topologie soll eine sichere, hochverfügbare und standardisierte Kommunikation zwischen dem lokalen Rechenzentrum, den Zweigstellen und den Cloud-Ressourcen etabliert werden.